\documentclass[11pt]{article}

\usepackage[margin=1in]{geometry}
\usepackage{amsmath,amssymb,amsthm,mathtools}
\usepackage{hyperref}

\newtheorem{theorem}{Theorem}
\newtheorem{proposition}{Proposition}
\newtheorem{lemma}{Lemma}
\newtheorem{corollary}{Corollary}
\newtheorem{remark}{Remark}

\newcommand{\R}{\mathbb R}
\newcommand{\ip}[2]{\left\langle #1,#2\right\rangle}
\newcommand{\norm}[1]{\left\lVert #1\right\rVert}
\newcommand{\dd}{\mathrm d}

\title{Residual Bounds and a Computable Self-Concordant Level Set}
\author{}
\date{}

\begin{document}
\maketitle

\section{Setup}

Let $\phi$ be a strictly convex $C^3$ barrier on a convex domain $\Omega\subseteq\R^n$. Define
\[
    g(x):=\nabla\phi(x),
    \qquad
    H(x):=\nabla^2\phi(x).
\]
The Hessian metric induces the primal and dual local norms
\[
    \norm{u}_x^2:=u^\top H(x)u,
    \qquad
    \norm{w}_{x,*}^2:=w^\top H(x)^{-1}w.
\]
Let $d_R$ denote the corresponding Riemannian geodesic distance.

For $x,y\in\Omega$, define the Jeffreys divergence
\[
    V(x,y):=\ip{x-y}{g(x)-g(y)}.
\]
We also define the endpoint primal-dual residual
\[
    d(x,y)
    :=
    \norm{x-y}_x^2
    +
    \norm{g(x)-g(y)}_{x,*}^2.
\]

Throughout the main self-concordant estimates we use the normalization
\[
    |D^3\phi(x)[a,b,b]|
    \le
    2\norm{a}_x\norm{b}_x^2.
\]
Equivalently, the Hessian metric satisfies the exponential stability estimate
\[
    e^{-2d_R(x,z)}H(x)
    \preceq
    H(z)
    \preceq
    e^{2d_R(x,z)}H(x).
\]

We will also use the standard geodesic coercivity inequality
\[
    d_R(x,y)^2\le V(x,y).
\]

\section{Two elementary consequences of exponential stability}

\begin{lemma}[Endpoint residual energy bound]
\label{lem:Ebound}
For all $x,y\in\Omega$,
\[
    d(x,y)
    \le
    2\left(e^{d_R(x,y)}-1\right)^2.
\]
Consequently, if $d_R(x,y)^2\le V(x,y)$, then
\[
    d(x,y)
    \le
    2\left(e^{\sqrt{V(x,y)}}-1\right)^2.
\]
Thus any geodesic step with kinetic energy $E$ satisfying
\[
    E\le d(x,y)
\]
obeys
\[
    E
    \le
    2\left(e^{\sqrt{V(x,y)}}-1\right)^2.
\]
\end{lemma}

\begin{proof}
Let $R=d_R(x,y)$ and let $\gamma:[0,R]\to\Omega$ be a unit-speed geodesic with
$\gamma(0)=x$ and $\gamma(R)=y$. Then
\[
    y-x=\int_0^R \gamma'(s)\,\dd s.
\]
Exponential stability gives
\[
    H(x)\preceq e^{2s}H(\gamma(s)),
\]
because $d_R(x,\gamma(s))=s$. Hence
\[
    \norm{\gamma'(s)}_x
    \le
    e^s\norm{\gamma'(s)}_{\gamma(s)}
    =e^s.
\]
Therefore
\[
    \norm{x-y}_x
    \le
    \int_0^R e^s\,\dd s
    =
    e^R-1.
\]

For the dual term,
\[
    g(y)-g(x)=\int_0^R H(\gamma(s))\gamma'(s)\,\dd s.
\]
By duality,
\[
\begin{aligned}
    \norm{g(y)-g(x)}_{x,*}
    &=
    \sup_{\norm{w}_x\le1}
    \int_0^R \ip{w}{H(\gamma(s))\gamma'(s)}\,\dd s  \\
    &\le
    \sup_{\norm{w}_x\le1}
    \int_0^R
    \norm{w}_{\gamma(s)}\norm{\gamma'(s)}_{\gamma(s)}\,\dd s.
\end{aligned}
\]
Since $H(\gamma(s))\preceq e^{2s}H(x)$, we have
\[
    \norm{w}_{\gamma(s)}\le e^s\norm{w}_x.
\]
The geodesic has unit speed, so
\[
    \norm{g(y)-g(x)}_{x,*}
    \le
    \int_0^R e^s\,\dd s
    =
    e^R-1.
\]
Squaring the primal and dual bounds and adding gives
\[
    d(x,y)\le 2(e^R-1)^2.
\]
The $V$-based bound follows from $R=d_R(x,y)\le\sqrt{V(x,y)}$.
\end{proof}

\begin{lemma}[Asymmetry residual bound]
\label{lem:Ubound}
Define
\[
    U(x,y):=(x-y)-H(x)^{-1}(g(x)-g(y)).
\]
Then
\[
    \norm{U(x,y)}_x
    \le
    2\left(\cosh(d_R(x,y))-1\right).
\]
Consequently, if $d_R(x,y)^2\le V(x,y)$, then
\[
    \norm{U(x,y)}_x
    \le
    2\left(\cosh\sqrt{V(x,y)}-1\right).
\]
\end{lemma}

\begin{proof}
Let $R=d_R(x,y)$, and let $\gamma:[0,R]\to\Omega$ be a unit-speed geodesic with
$\gamma(0)=x$ and $\gamma(R)=y$. Then
\[
    x-y=-\int_0^R \gamma'(s)\,\dd s,
\]
and
\[
    g(x)-g(y)=-\int_0^R H(\gamma(s))\gamma'(s)\,\dd s.
\]
Thus
\[
\begin{aligned}
    U(x,y)
    &=(x-y)-H(x)^{-1}(g(x)-g(y)) \\
    &=
    \int_0^R
    \left[H(x)^{-1}H(\gamma(s))-I\right]\gamma'(s)\,\dd s.
\end{aligned}
\]
Set
\[
    A_s:=H(x)^{-1/2}H(\gamma(s))H(x)^{-1/2},
    \qquad
    z_s:=H(x)^{1/2}\gamma'(s).
\]
Exponential stability gives
\[
    e^{-2s}I\preceq A_s\preceq e^{2s}I.
\]
Since $\gamma$ has unit speed,
\[
    z_s^\top A_s z_s
    =
    \gamma'(s)^\top H(\gamma(s))\gamma'(s)
    =1.
\]
The $x$-norm of the integrand equals
\[
    \left\|
    \left[H(x)^{-1}H(\gamma(s))-I\right]\gamma'(s)
    \right\|_x
    =
    \norm{(A_s-I)z_s}_2.
\]
Diagonalizing $A_s$, and using $z_s^\top A_s z_s=1$, gives
\[
    \norm{(A_s-I)z_s}_2^2
    \le
    \max_{\lambda\in[e^{-2s},e^{2s}]}
    \frac{(\lambda-1)^2}{\lambda}.
\]
The maximum is attained at an endpoint and equals
\[
    e^{2s}+e^{-2s}-2
    =
    (e^s-e^{-s})^2.
\]
Therefore
\[
    \left\|
    \left[H(x)^{-1}H(\gamma(s))-I\right]\gamma'(s)
    \right\|_x
    \le
    e^s-e^{-s}.
\]
Integrating yields
\[
    \norm{U(x,y)}_x
    \le
    \int_0^R (e^s-e^{-s})\,\dd s
    =
    e^R+e^{-R}-2
    =
    2(\cosh R-1).
\]
The $V$-based bound follows from $R\le\sqrt{V(x,y)}$ and monotonicity of $\cosh R-1$ for $R\ge0$.
\end{proof}

\section{The target-dependent one-step estimate}

Let $x(t)$, $t\in[0,1]$, be a unit-time Hessian geodesic with
\[
    x(0)=x,
    \qquad
    x(1)=T(x),
\]
and constant kinetic energy
\[
    E:=\norm{\dot x(t)}_{x(t)}^2.
\]
Fix a target point $y$ and define
\[
    v(t):=V(x(t),y),
    \qquad
    v_0:=v(0)=V(x,y).
\]

\begin{lemma}[Second-variation identity]
\label{lem:secondvariation}
Along $x(t)$,
\[
    v''(t)=2E-\ip{\ddot x(t)}{U(x(t),y)}_{x(t)},
\]
where $\ip{a}{b}_z:=a^\top H(z)b$.
Consequently,
\[
    v(1)
    =
    v_0+v'(0)+E
    -
    \int_0^1(1-t)\ip{\ddot x(t)}{U(x(t),y)}_{x(t)}\,\dd t.
\]
\end{lemma}

\begin{proof}
Differentiating
\[
    v(t)=\ip{x(t)-y}{g(x(t))-g(y)}
\]
gives
\[
    v'(t)
    =
    \ip{\dot x(t)}{g(x(t))-g(y)}
    +
    \ip{x(t)-y}{H(x(t))\dot x(t)}.
\]
Differentiating again yields
\[
\begin{aligned}
    v''(t)
    &=
    \ip{\ddot x(t)}{g(x(t))-g(y)}
    +2\dot x(t)^\top H(x(t))\dot x(t) \\
    &\quad
    +D^3\phi(x(t))[\dot x(t),\dot x(t),x(t)-y]
    +\ip{x(t)-y}{H(x(t))\ddot x(t)}.
\end{aligned}
\]
The Hessian geodesic equation implies that for every vector $w$,
\[
    \ip{\ddot x(t)}{w}_{x(t)}
    =
    -\frac12D^3\phi(x(t))[\dot x(t),\dot x(t),w].
\]
Using this with $w=x(t)-y$ and using
$\dot x(t)^\top H(x(t))\dot x(t)=E$, we obtain
\[
\begin{aligned}
    v''(t)
    &=
    2E
    +\ip{\ddot x(t)}{g(x(t))-g(y)}
    -\ip{\ddot x(t)}{x(t)-y}_{x(t)} \\
    &=
    2E
    -\ip{\ddot x(t)}{(x(t)-y)-H(x(t))^{-1}(g(x(t))-g(y))}_{x(t)}.
\end{aligned}
\]
This is the claimed identity. Taylor's formula with integral remainder gives
\[
    v(1)=v_0+v'(0)+\int_0^1(1-t)v''(t)\,\dd t,
\]
and $\int_0^1(1-t)2E\,\dd t=E$.
\end{proof}

\begin{lemma}[Acceleration bound]
\label{lem:acceleration}
For a standard self-concordant Hessian geodesic,
\[
    \norm{\ddot x(t)}_{x(t)}\le E.
\]
\end{lemma}

\begin{proof}
By the geodesic equation,
\[
    \ip{\ddot x(t)}{w}_{x(t)}
    =
    -\frac12D^3\phi(x(t))[\dot x(t),\dot x(t),w].
\]
The self-concordance inequality gives
\[
    |D^3\phi(x(t))[w,\dot x(t),\dot x(t)]|
    \le
    2\norm{w}_{x(t)}\norm{\dot x(t)}_{x(t)}^2
    =
    2E\norm{w}_{x(t)}.
\]
Since the cubic tensor is symmetric,
\[
    |\ip{\ddot x(t)}{w}_{x(t)}|
    \le
    E\norm{w}_{x(t)}.
\]
Taking the supremum over $\norm{w}_{x(t)}\le1$ gives
\[
    \norm{\ddot x(t)}_{x(t)}\le E.
\]
\end{proof}

\begin{theorem}[Symbolic one-step bound]
\label{thm:symbolic}
Assume the descent condition
\[
    v'(0)
    \le
    -(E+v_0),
\]
the step-energy compatibility condition
\[
    E\le d(x,y),
\]
and the path sublevel condition
\[
    V(x(t),y)
    \le
    v_0
    \qquad
    \forall t\in[0,1].
\]
Then
\[
    V(T(x),y)
    \le
    f(v_0),
\]
where
\[
    f(v)
    :=
    2\left(e^{\sqrt v}-1\right)^2
    \left(\cosh\sqrt v-1\right).
\]
In particular,
\[
    f(v)=v^2+O(v^{5/2})
    \qquad
    \text{as }v\downarrow0.
\]
\end{theorem}

\begin{proof}
By Lemma~\ref{lem:secondvariation},
\[
    v(1)
    =
    v_0+v'(0)+E
    -
    \int_0^1(1-t)\ip{\ddot x(t)}{U(x(t),y)}_{x(t)}\,\dd t.
\]
The descent condition implies
\[
    v_0+v'(0)+E\le0.
\]
Therefore
\[
    v(1)
    \le
    \int_0^1(1-t)
    |\ip{\ddot x(t)}{U(x(t),y)}_{x(t)}|\,\dd t.
\]
By Lemma~\ref{lem:acceleration},
\[
    |\ip{\ddot x(t)}{U(x(t),y)}_{x(t)}|
    \le
    E\norm{U(x(t),y)}_{x(t)}.
\]
The path sublevel condition and Lemma~\ref{lem:Ubound} give
\[
    \norm{U(x(t),y)}_{x(t)}
    \le
    2(\cosh\sqrt{v_0}-1)
    \qquad
    \forall t\in[0,1].
\]
Therefore
\[
\begin{aligned}
    v(1)
    &\le
    \int_0^1(1-t)
    E\cdot 2(\cosh\sqrt{v_0}-1)\,\dd t \\
    &=
    E(\cosh\sqrt{v_0}-1).
\end{aligned}
\]
Finally, by the step-energy compatibility condition and Lemma~\ref{lem:Ebound},
\[
    E\le d(x,y)
    \le
    2(e^{\sqrt{v_0}}-1)^2.
\]
Substituting this into the previous display yields
\[
    v(1)
    \le
    2(e^{\sqrt{v_0}}-1)^2(\cosh\sqrt{v_0}-1)
    =f(v_0).
\]
The expansion follows from
\[
    (e^{\sqrt v}-1)^2=v+O(v^{3/2}),
    \qquad
    \cosh\sqrt v-1=\frac12v+O(v^2).
\]
\end{proof}

\section{A concrete numerical level set}

\begin{corollary}[Concrete contraction level]
Let
\[
    f(v)
    =
    2\left(e^{\sqrt v}-1\right)^2
    \left(\cosh\sqrt v-1\right).
\]
The equation
\[
    f(v)=v
\]
has a positive root
\[
    v_\star
    \approx
    0.4671709662.
\]
Moreover,
\[
    f(v)\le v
    \qquad
    \text{for }0\le v\le v_\star.
\]
Consequently, under the hypotheses of Theorem~\ref{thm:symbolic}, if
\[
    V(x,y)\le v_\star
\]
and the path sublevel condition holds, then
\[
    V(T(x),y)\le V(x,y).
\]
Thus $\mathcal L_{v_\star}(y):=\{x:V(x,y)\le v_\star\}$ is a computable contraction level set for this estimate.
\end{corollary}

\begin{proof}
This is a one-dimensional numerical verification.  Solving
\[
    2(e^{\sqrt v}-1)^2(\cosh\sqrt v-1)=v
\]
gives
\[
    v_\star\approx0.4671709662.
\]
A direct check of the scalar function on the interval $[0,v_\star]$ gives
$f(v)\le v$.  Since Theorem~\ref{thm:symbolic} gives
$V(T(x),y)\le f(V(x,y))$, the result follows.
\end{proof}

\begin{corollary}[Concrete quadratic-rate bound]
On the smaller level set $V(x,y)\le 1/4$, the same scalar estimate gives
\[
    V(T(x),y)
    \le
    1.719\,V(x,y)^2.
\]
More precisely,
\[
    \sup_{0<v\le1/4}\frac{f(v)}{v^2}
    \approx
    1.718720647.
\]
Thus this estimate gives standard quadratic convergence on $\mathcal L_{1/4}(y)$, with explicit quadratic constant $1.719$.
\end{corollary}

\begin{proof}
The scalar ratio $f(v)/v^2$ is increasing on $0<v\le1/4$ for this function.  Hence
\[
    \sup_{0<v\le1/4}\frac{f(v)}{v^2}
    =
    \frac{f(1/4)}{(1/4)^2}
    \approx
    1.718720647.
\]
Combining this with Theorem~\ref{thm:symbolic} gives the claim.
\end{proof}

\begin{remark}[Normalized quadratic contraction]
The scalar estimate $f(v)$ satisfies
\[
    f(v)=v^2+O(v^{5/2}).
\]
For small positive $v$, the higher-order correction is positive, so this particular bound does not certify the normalized estimate $f(v)\le v^2$ on any interval $0<v\le\delta$.  It does certify the usual quadratic rate $f(v)\le C_\delta v^2$ on every sufficiently small level set.
\end{remark}

\end{document}
